\section{Introduction}

Let $G$ be a finite simple undirected graph with $N$ vertices $n_1,\dots,n_N$
and adjacency matrix $A$ (defined in \cref{sec:setting}). The lecture board
asks two questions about it, and the assignment adds a third
\cite{whiteboard1,whiteboard2,userstatements}.
\begin{enumerate}
\item[(Q1)] What does the entry $(A^k)_{ij}$ represent? The mark above $A$ on
  the board is read as an exponent $k$ \cite{userstatements}.
  \emph{Answer:} it is the number of walks of length $k$ from $n_i$ to $n_j$
  (\cref{thm:counting}); it counts walks, not paths.
\item[(Q2)] How can one tell from $A$ whether $G$ is connected?
  \emph{Answer:} $G$ is connected if and only if every entry off the
  diagonal of $B=A+A^2+\dots+A^{N-1}$ is positive (\cref{thm:sum}).
\item[(Q3)] What is a path between $n_i$ and $n_j$? \emph{Answer:} a walk
  from $n_i$ to $n_j$ in which no vertex occurs twice
  (\cref{def:pathbetween}).
\end{enumerate}
The answers come in the order Q1, Q3, Q2, because the test for Q2 needs the path. A full-text search of the course textbook
\cite{petrosyan2015graphs} finds neither the walk-counting theorem nor a matrix test for
connectivity, so both are taken from other sources, and every proof is written out in
full; ``derived'' marks this paper's own proof of a result stated in a cited source.
Textbook page numbers are printed ones, one lower than the PDF's.

Plain-words paragraphs precede the formal definitions, theorems, lemmas, and the worked example. They tell one story. A pond has \emph{islands}, and some pairs of islands have a \emph{bridge} between them (vertices and edges); two islands joined by a bridge are \emph{next door}. A frog makes \emph{trips} of \emph{hops}, each hop crossing one bridge, and it may hop back or land twice on an island (a walk); a \emph{fresh-island trip} lands on no island twice (a path), and the pond is \emph{in one piece} when the frog can get from any island to any other (connected). A square \emph{table} with a one in each cell whose two islands are next door, and a zero elsewhere, is the adjacency matrix.
